\section{Research Approach (Methodology)}
The empirical study will employ an experimental systems methodology based on controlled, repeatable cluster deployments. A Kubernetes-in-Docker (Kind) environment will provide the primary testbed to ensure reproducibility and rapid scenario iteration, while still exposing realistic orchestration behavior for service discovery, policy updates, and workload scaling. Two principal architecture variants will be instantiated: a sidecar-based baseline and a sidecarless configuration using eBPF-enabled service networking with Cilium and Istio Ambient Mesh components, including ztunnel where applicable \cite{kindProject,istioAmbient,ciliumArchitecture}.

Load generation will use Fortio and wrk2 to capture both closed-loop and open-loop traffic regimes, enabling robust analysis of queue buildup and tail amplification under varying offered load \cite{fortio,wrk2}. Experimental factors will include tenant count, request concurrency, burstiness, resource saturation level, and control-plane update frequency. The primary response variable will be p99 latency, supplemented by p95, throughput, loss, and jitter to support interpretation. The design will prioritize repeated trials and confidence-aware reporting to reduce sensitivity to run-to-run variance.

Observation and diagnosis will combine application-level metrics with kernel-level evidence. bpftrace will be used for targeted tracepoint and kprobe-level inspection of scheduling and networking events, while Cilium Hubble will provide flow-level context for policy and path decisions \cite{bpftrace,ciliumHubble}. This dual-observability approach is intended to bridge macro-level latency behavior with micro-level causal signals. Statistical analysis will compare distributions across architectural modes and disturbance regimes, with particular attention to tail leakage patterns and the conditions under which sidecarless efficiency gains degrade into fairness trade-offs.
